\documentclass[11pt]{article}
\usepackage[margin=1in]{geometry}
\usepackage{amsmath,amssymb}
\usepackage{graphicx}
\usepackage{booktabs}
\usepackage{hyperref}
\usepackage{siunitx}
\hypersetup{colorlinks=true,linkcolor=blue,urlcolor=blue}
\title{Replication Report:\\
\emph{Magnetic octupoles as the order parameter for unconventional
antiferromagnetism}\\
\large Bhowal \& Spaldin, arXiv:2212.03756 (v.\ arXiv:2205.09500), Phys.\ Rev.\ X \textbf{14}, 011019 (2024)}
\author{Automated replication (minimal tight-binding + symmetry model)}
\date{2026-07-19}
\begin{document}
\maketitle

\begin{abstract}
We replicate the tractable core of Bhowal \& Spaldin's proposal that
\emph{ferroically ordered magnetic octupoles} are the natural order parameter of
centrosymmetric, time-reversal-broken antiferromagnets with non-relativistic
spin splitting (NRSS)---the ``altermagnet'' class---using rutile MnF$_2$ as the
prototype. Full density-functional theory (Elk/VASP LDA+SOC+$U$, atomic-site
multipole decomposition, and magnetic Compton profiles) is \textbf{out of scope}
for this replication. Instead we re-implement the paper's own minimal four-band
(eight-band with spin) tight-binding (TB) model [Eqs.~(2)--(6), Table~I] and its
reciprocal-space octupole symmetry analysis, and verify five/ten
machine-checkable claims. All checks pass. The analytic spin-splitting formula
Eq.~(6) reproduces the full $8\times8$ diagonalization to machine precision; the
peak $\Gamma\!\to\!M$ splitting is $41.2$~meV (paper Fig.~3d scale
$\sim\!30$--$40$~meV); the splitting is even in $\mathbf{k}$, exhibits $d$-wave
$C_4$ sign reversal, vanishes on the $k_x{=}0/k_y{=}0$ nodal lines, and is
controlled entirely by the two inter-sublattice hoppings $t_3,t_4$---exactly the
``structure--spin correlation'' that the octupole encodes.
\textbf{Verdict: strong replication of the model-level claims.
Coverage 6/10, Agreement 9/10.}
\end{abstract}

\section{The paper in brief}
The authors argue that in a compensated ($\mathbf{M}=0$), centrosymmetric AFM
whose bands are spin-split without spin--orbit coupling (SOC), the lowest-order
\emph{ferroically} ordered magnetic multipole allowed by symmetry is the rank-3
\emph{magnetic octupole}
\begin{equation}
O_{ijk}=\int r_i r_j\,\mu_k(\mathbf{r})\,d^3r ,
\end{equation}
the third term in the multipole expansion of the magnetic interaction energy
[Eq.~(1)]. The dipole term vanishes (no net moment), and the magnetoelectric
term $M_{ij}=\int r_i\mu_j\,d^3r$ is forbidden by inversion, leaving the octupole
as the \emph{natural order parameter}. For MnF$_2$ (space group
$P4_2/mnm$; magnetic space group $P4_2'/mnm'$) the ferro-type component is
$O_{32}^{-}$ with real-space form $xy\,m_z$ and reciprocal-space form
$k_xk_y m_z$, which \emph{directly explains} the $d$-wave NRSS along $\Gamma\to M$.
The octupole picture further (i) accounts for the known piezomagnetic effect
(rank-3 tensor $\Lambda_{xyz}{=}\Lambda_{yxz}{\ne}\Lambda_{zxy}$), (ii) predicts a
new \emph{anti}-piezomagnetic effect from the antiferro $O_{30}$ octupole, and
(iii) predicts a non-zero magnetic Compton profile as a direct experimental probe.

\section{What we replicated (scope)}
\textbf{In scope (this report):} the minimal TB model, its analytic eigenvalues
and spin splitting, the reciprocal-space octupole form factor, the $d$-wave
symmetry/sign-reversal structure, the role of microscopic hoppings, and the
symmetry of the piezomagnetic tensor implied by the ferro $O_{32}^-$ octupole.

\noindent\textbf{Out of scope (not attempted; require heavy first-principles
machinery / proprietary Elk extensions):} the LDA+SOC+$U$ electronic structure,
the atomic-site tensor-moment decomposition of the DFT density matrix
$\rho_{lm,l'm'}$ (Fig.~2), the band-decomposed charge/spin densities (Fig.~2e--h),
the shear-strain DFT relaxations giving the piezo/anti-piezomagnetic magnitudes
(Fig.~4), and the Elk magnetic Compton profiles (Fig.~5). These are flagged as
out-of-scope, not as failures.

\section{Model}
Basis $\{$Mn1-$d_{xz}$, Mn1-$d_{yz}$, Mn2-$d_{xz}$, Mn2-$d_{yz}\}$;
$\vec\Sigma$ (sublattice) and $\vec\sigma$ (orbital) Pauli matrices. The
spinless Hamiltonian [Eq.~(2)] is
\begin{equation}
H_t=\alpha(\mathbf{k})\,\mathbb{I}+\beta(\mathbf{k})\,\Sigma_z\!\otimes\!\sigma_x
+\gamma(\mathbf{k})\,\Sigma_x\!\otimes\!\sigma_0
+\delta(\mathbf{k})\,\Sigma_x\!\otimes\!\sigma_x ,
\end{equation}
with $\alpha=\varepsilon_1+2t_1\cos k_zc$, $\beta=\varepsilon_2+2t_2\cos k_zc$,
$\gamma=8t_3\cos\frac{k_xa}{2}\cos\frac{k_ya}{2}\cos\frac{k_zc}{2}$,
$\delta=-8t_4\sin\frac{k_xa}{2}\sin\frac{k_ya}{2}\cos\frac{k_zc}{2}$ [Eq.~(3)].
Adding AFM exchange [Eq.~(5)] $H=S_0\otimes H_t+J\,S_z\otimes(\Sigma_z\otimes\sigma_0)$
gives the top-valence spin splitting [Eq.~(6)]
\begin{equation}
\Delta E_s=\sqrt{(J{+}\beta)^2+(\delta{-}\gamma)^2}-\sqrt{(J{+}\beta)^2+(\delta{+}\gamma)^2}
\;\approx\;\frac{32}{\epsilon}\,t_3 t_4\,\sin(k_xa)\sin(k_ya),
\quad \epsilon=J+\varepsilon_2+2t_2 .
\end{equation}
We use the verbatim NMTO parameters of Table~I (in Ry):
$t_1{=}0.0036,\ t_2{=}{-}0.0038,\ t_3{=}0.0040,\ t_4{=}0.0034,\
\varepsilon_1{=}{-}0.1385,\ \varepsilon_2{=}{-}0.0339$, and $2J\approx5$~eV.
Two independent implementations were written: \texttt{tb\_mnf2.py} (eV units,
realistic $a{=}4.8734$, $c{=}3.3099$~\AA) and \texttt{model.py}/\texttt{run\_checks.py}
(reduced units); they agree.

\section{Results}
\begin{table}[h]\centering
\caption{Machine-checkable claims (canonical suite \texttt{verify\_claims.py}; all
values from \texttt{work/verify\_results.json}). 10/10 pass.}
\begin{tabular}{clc}
\toprule
\# & Claim & Result \\
\midrule
1 & Analytic Eq.~(6) $=$ full $8\times8$ diag (2000 random $\mathbf{k}$) & max err $6.2\times10^{-15}$~eV \checkmark\\
2 & Approx.\ $\frac{32}{\epsilon}t_3t_4\sin k_x\sin k_y$ vs exact ($\Gamma\!\to\!M$) & max rel err $2.5\%$ \checkmark\\
3 & $\Delta E_s(\mathbf{k})=\Delta E_s(-\mathbf{k})$ (even, not Rashba) & max asym $0$ \checkmark\\
4 & $C_4$ sign flip $(k_x,k_y)\!\to\!(k_x,-k_y)$, $|\cdot|$ preserved & 20/20 flip \checkmark\\
5 & Nodal lines: $\Delta E_s{=}0$ on $k_x{=}0$ or $k_y{=}0$ & max $0$ on nodes \checkmark\\
6 & Splitting requires both $t_3,t_4$ (set either $\to0\Rightarrow0$) & $0$ eV \checkmark\\
7 & Peak $|\Delta E_s|$ on Fig.~3d scale & $41.2$~meV \checkmark\\
8 & Tracks $k_xk_y$ ($B_{1g}$/$xy$) octupole form factor & Pearson $0.99999$ \checkmark\\
9 & Sign reverses for modified structure / opposite AFM domain & mag.\ preserved \checkmark\\
10& Ferro $xy\,m_z$ octupole $\Rightarrow$ $O_{z,xy}{=}\!\int x^2y^2{\ne}0$; $O_{z,xx}{=}0$ & \checkmark\\
\bottomrule
\end{tabular}
\end{table}

\begin{figure}[h]\centering
\includegraphics[width=0.98\textwidth]{../work/fig_spin_splitting.png}
\caption{\textbf{Left:} reproduction of paper Fig.~3(d). The exact Eq.~(6)
(blue), the analytic $d$-wave approximation (red dashed) and the brute-force
$8\times8$ diagonalization (green dots) overlay; the curve is symmetric about
$\Gamma$ and \emph{reverses sign} between $(0.5,0.5,0)$ and $(0.5,-0.5,0)$,
the hallmark $d$-wave NRSS. \textbf{Right:} $\Delta E_s(k_x,k_y)$ in the
$k_z{=}0$ plane---a clean four-lobe $d$-wave ($B_{1g}$) pattern with nodes on the
axes, i.e.\ the reciprocal-space image of the $O_{32}^-=k_xk_y m_z$ octupole.}
\end{figure}

\section{Quantitative comparison with the paper}
\begin{itemize}
\item \textbf{Peak spin splitting} $\Gamma\!\to\!M$: replication $41.2$~meV;
paper Fig.~3(d) DFT peak $\approx-30$ to $-40$~meV. Same order and $k$-location
($k\approx0.5\,\Gamma M$). Sign is convention-dependent (up/down labelling).
\item \textbf{Analytic vs.\ exact} (Eq.~6 last equality vs.\ first): mean
relative error $1.6\%$, max $2.5\%$ along $\Gamma\!\to\!M$; the paper states the
two ``agree reasonably''---consistent.
\item \textbf{Eq.~(6) vs.\ full diagonalization:} exact to $6\times10^{-15}$~eV,
confirming the analytic reduction the authors perform.
\item \textbf{Symmetry:} even in $\mathbf{k}$, $d$-wave $C_4$ sign reversal, and
axis nodes are reproduced exactly, matching the $k_xk_y m_z$ reciprocal
representation and the reported $C_4$ reversal (Fig.~3b,c).
\item \textbf{Microscopic origin:} splitting $\propto t_3 t_4$ (inter-sublattice
intra/inter-orbital hoppings); zeroing either kills it---the paper's central
mechanistic statement.
\end{itemize}

\section{Conclusion}
The model-level physics of the paper is fully and independently reproducible from
the equations and parameters as published. The octupole $\leftrightarrow$ NRSS
$\leftrightarrow$ piezomagnetism logical chain holds at the symmetry level. The
DFT-derived numbers (multipole magnitudes in $\mu_B$, Compton profile
amplitudes, strain-induced moments) are out of scope and unverified here.
\textbf{Verdict: strong replication (model tier). Coverage 6/10
(model + symmetry reproduced; DFT/Compton/strain not attempted).
Agreement 9/10 (every attempted claim matches, quantitatively where numbers
exist).}
\end{document}
